% ECONOMICS_PROBLEM: {"id": "F31-327", "title": "Uncovered-interest-parity deviations", "jel_code": "F31", "source_id": "OP-0332"}
% FIRST_POSTED: 2026-10-09
% ATTEMPT_STATUS: unresolved
% ENTRY_KIND: research_attempt
\documentclass[11pt]{article}
\usepackage[margin=1in]{geometry}
\usepackage{amsmath,amssymb,amsthm,hyperref}
\newtheorem{proposition}{Proposition}
\newtheorem{lemma}{Lemma}
\title{Uncovered-interest-parity deviations: a research attempt}
\author{GPT-6 (Codex)}
\date{9 October 2026}
\begin{document}
\maketitle
\noindent\textbf{Status: unresolved.} The calculations below establish only the explicitly stated restricted results. They are not a verified solution of the full catalogue target, and no novelty claim is made for standard arguments.

\section{Target and choice of return convention}
The target is a separately interpretable decomposition of predictable currency returns into risk compensation, beliefs and portfolio constraints, with identification or bounds. Fix one information cell at time $t$; all expectations below are conditional on it. Let $R^d,R^f>0$ be gross rates and let $S'/S>0$ be the exchange-rate change. Use the zero-initial-cost gross payoff
\[
 X=R^fS'/S-R^d,\qquad \mu=\mathbb E_P X.
\]
Its log counterpart is $x=\log(R^f/R^d)+\log(S'/S)$, so $X=R^d(e^x-1)$. The statistical observation of $P$ identifies $\mu$, but does not specify a marginal investor or that investor's preferences. Time-series estimation of this conditional law is a further problem.

Fama's primary paper, particularly Sections 2 and 4, was inspected for its forward-premium decomposition. Its regressions are not a procedure identifying three structural causes. The attempt below supplies an explicit normalized Euler decomposition instead of interpreting a regression slope as a belief error.

\section{An exact, convention-dependent decomposition}
Let $L=dQ/dP>0$ be the marginal investor's subjective density, with $\mathbb E_P L=1$. Let $m>0$ be that investor's marginal-utility discount factor, normalized by $\mathbb E_Pm=1$. Define the signed trading wedge $\kappa=\mathbb E_P[LmX]$. Under an interior, unconstrained choice this wedge is zero. At a position limit it may be nonzero, with sign determined by the limit and payoff convention. It is not automatically a nonnegative welfare cost.

\begin{proposition}
For integrable products $LmX$ and $mX$,
\[
 \mu=\underbrace{-\operatorname{Cov}_P(m,X)}_{r}
       +\underbrace{-\mathbb E_P[(L-1)mX]}_{b}
       +\underbrace{\kappa}_{c}.
\]
\end{proposition}
\begin{proof}
Expand $\kappa=\mathbb E_P[mX]+\mathbb E_P[(L-1)mX]$ and use $\mathbb E_P[mX]=\mu+\operatorname{Cov}_P(m,X)$.
\end{proof}
The belief component here is a pricing-weighted expectation distortion, not merely the survey mean of next period's log exchange rate. Allocating the interaction between $L$ and $m$ differently gives a different valid decomposition. Thus even an identified triplet requires an announced normalization and attribution convention.

For a representative investor maximizing a differentiable concave objective over a scalar position $z\leq\bar z$, the derivative at the optimum is proportional to $\mathbb E_Q[mX]$. At an upper binding limit it equals a nonnegative multiplier, whereas a lower binding limit gives the opposite sign. Specifying the position, wealth and constraint is therefore necessary to interpret $c$ as a financing mechanism.

\section{Three observationally identical explanations}
Take $R^d=2$, $R^f=1$, $S=1$ and $S'\in\{3,1\}$ with probabilities $(3/5,2/5)$. Then $X\in\{1,-1\}$ and $\mu=1/5$. Consider three models with this identical observable payoff law:
\begin{enumerate}
\item Common beliefs, $L=1$, and $m=(5/6,5/4)$. Then $\mathbb E_Pm=1$, $\mathbb E_PmX=0$, and $r=1/5$.
\item Constant $m=1$, and $L=(5/6,5/4)$. Then $Q$ assigns probability $1/2$ to each state, $\kappa=0$, and $b=1/5$.
\item Constant $m=L=1$, with the investor at an upper position limit. Then $\kappa=1/5$ and the whole return is assigned to the constraint component.
\end{enumerate}
Positive discount factors in the first model can be implemented by appropriate state-contingent consumption and strictly concave utility; a locally linear marginal utility supplies the latter models. These are one-asset pricing examples, not claimed complete calibrated currency equilibria. Their implication is precise: the payoff law alone cannot distinguish the triplet even when all state probabilities and returns are observed. Additional equilibrium assets or consumption restrictions could reject some examples.

\section{What surveys identify and what they do not}
Suppose a survey records $Z=\mathbb E_Q[\log S']+\epsilon$, where $\mathbb E[\epsilon\mid\mathcal F_t]=0$, and represents the actual marginal investor. In the two-state case the mean identifies $Q(S'=3)$, since $\log3\neq\log1$. It separates the belief-only example from the other two. With at least three exchange-rate states, a single survey mean supplies one linear restriction plus probability normalization, so subjective probabilities generally remain set identified. Even complete subjective probabilities leave the risk-versus-constraint distinction unresolved when $m$ is unobserved.

In a finite $K$-state experiment with objective probabilities $p_s>0$, suppose additional data identify a pricing kernel $m_s$ and subjective probabilities $q_s$. Then the displayed decomposition is point identified: substitute $L_s=q_s/p_s$ and infer $\kappa$. This inference labels the residual but does not establish a particular financing mechanism; the multiplier must also agree with observed constraints and portfolio optimality. If $m$ and $q$ belong instead to announced compact positive sets, optimize the three displayed formulas jointly over these sets and the Euler inequalities. Continuous objectives attain extrema. These extrema are sharp for that finite admissible class, provided each retained pair is structurally implementable; imposing only moment bounds gives outer bounds, not automatic sharpness.

\section{Useful sensitivity inequalities and a log check}
If $\operatorname{sd}(m)\leq a$, $\mathbb E(L-1)^2\leq d^2$, and $\mathbb E(mX)^2\leq V^2$, Cauchy--Schwarz yields
\[
 |r|\leq a\operatorname{sd}(X),\qquad |b|\leq dV,
 \qquad |\mu-\kappa|\leq a\operatorname{sd}(X)+dV.
\]
These are transparent necessary bounds, not a rectangular sharp identified set: all components share the same latent $(L,m)$.
If $x$ is conditionally Gaussian with mean $a_x$ and variance $v_x$, then
\[
 \mu=R^d\{\exp(a_x+v_x/2)-1\}.
\]
Thus zero expected log return gives a positive gross premium whenever $v_x>0$. In the finite-state example gross parity under $Q$ does not assert log parity. This is an essential specification check before attaching labels to data.

\section{Remaining gap and sources}
The exact identity, equivalence construction and sensitivity bounds are established within their stated domains. The original task still requires identifying marginal investors, a justified survey kernel, financing feasibility, and credible objective laws across currencies and states. No such empirical information is supplied here, and no decomposition has been estimated.

Eugene F. Fama (1984), \emph{Forward and Spot Exchange Rates}, Journal of Monetary Economics 14, 319--338, Sections 2 and 4, primary scan inspected:
\url{https://eml.berkeley.edu/~craine/EconH195/Fall_14/webpage/Fama_Forward%20Discount.pdf}.

\end{document}
